\subsubsection{Training-only Feature Selection}
Both candidate selectors retain 20 original source fields. CART importance selection uses a decision tree with maximum depth 5 and minimum leaf size 4, retaining the 20 largest impurity importances. Recursive feature elimination (RFE) uses the same CART estimator, removes one field per step and refits until 20 remain. Imputation precedes selection, and the selected continuous fields are transformed after selection; ordinal category codes remain unchanged.

Both methods are supervised and therefore fitted only on training records. Every inner CV fit repeats imputation and selection, so its validation records cannot influence the retained fields. Stacking maintains a complete preparation pipeline for each base learner and refits selection inside its own out-of-fold construction. After tuning, the final pipeline refits selection using the complete 223-record training partition, before predicting the 96 held-out records.

Selected-feature CSVs include source names, selection flags and ranks, separately for each stacking base learner. Selection may vary across seeds and fits. Impurity importance can favour predictors with more available split points and does not establish a clinical effect. The four controlled preparation benchmarks in Section 8 evaluate selector and transformation changes while maintaining fixed model settings and test membership.
